\section{Registered Measurement Protocol}
\label{sec:protocol}

\paragraph{Data and split.}
We use NIH ChestX-ray14 \citep{wang2017chestxray}, whose findings are mined from radiology reports. A deterministic hash of patient identity assigns rows to disjoint train, validation, and test partitions. The registered manifest contains 18,212 training images and 5,343 test images from 1,659 test patients. The three cells pair LLaVA-1.5-7B with Effusion and Edema, and Qwen2.5-VL-7B with Effusion. All reported uncertainty clusters by patient. Prompts and exact run identifiers appear in Appendix~\ref{app:protocol}.

\paragraph{Models and consumed loci.}
For LLaVA-1.5-7B \citep{liu2024improved}, \texttt{vis.last} is the output of \texttt{model.vision\_tower.encoder.layers.22}; for Qwen2.5-VL-7B \citep{bai2025qwen25vl}, it is \texttt{model.visual.blocks.31}. Synthetic hook tests establish that each module is on the downstream path, that an $\alpha=0$ intervention is bitwise identical to the unmodified forward pass, and that a nonzero intervention changes the connector or merger and final logits. The probe and intervention therefore operate on the same final visual representation that downstream computation consumes.

\paragraph{Controlled decodability.}
For image $i$, mean pooling over the final-block visual tokens yields $\hvec_i\in\R^D$. A seed-0 Gaussian projection maps every model to 512 dimensions; the scaler is fitted on training rows only. We then fit logistic regression with $C=1$, maximum 2,000 iterations, and seed 0. The real-label AUROC is compared with 20 control AUROCs. Each control seed maps the same recurring input type---AP view indicator, sex indicator, and age decade, giving 39 strata---to a random label. Selectivity is
\begin{equation}
S = \operatorname{AUROC}(y,\hat y)-\frac{1}{20}\sum_{s=0}^{19}\operatorname{AUROC}(\tilde y_s,\hat y_s).
\label{eq:selectivity}
\end{equation}
We report the real AUROC, control distribution, and $S$ with 2,000 patient-cluster bootstrap resamples. Paired contrasts reuse identical ordered test rows and bootstrap draws.

\paragraph{Probe-normal intervention.}
Let $\hat\wvec_c$ be the unit normal of the capacity-matched probe for concept $c$. At every visual token $t$, the registered intervention is
\begin{equation}
\hvec'_{it}=\hvec_{it}+\alpha\lVert\hvec_{it}\rVert_2\hat\wvec_c.
\label{eq:intervention}
\end{equation}
This relative-token scale holds perturbation magnitude fixed as a fraction of the local activation norm. The endpoint is the paired change in mean $P(yes)$ from the common unsteered baseline. Original cells use concept doses $\{-1,-0.5,-0.25,-0.1,0,0.1,0.25,0.5,1\}$ on 200 held-out test images. Controls are evaluated at the six nonzero doses excluding $\pm0.1$: 20 isotropic random unit directions, one coordinate-permutation sham, and fixed unrelated clinical probe normals. Dose monotonicity over $[-0.5,0.5]$ is descriptive.

\paragraph{Original-cell decision rule.}
For each controlled nonzero dose, we orient the raw change by the sign registered for the concept direction. The primary effect is the maximum concept-consistent change over those doses. A cell clears the generic intervention rung only if this effect leaves the same-dose random 5th--95th percentile interval in the concept-consistent direction and exceeds the maximum absolute sham effect. The rule is fixed across the three cells.

\paragraph{Direction-specificity supplement.}
The original Qwen cell motivated a prospective independent-row test at its discovery-locked dose, $\alpha=+0.25$. We exclude all patients in the original 200-image cohort, select one row per remaining test patient by a label-blind hash, rank patients by a second label-blind hash, and take 400. The grid contains a common unsteered baseline, Effusion, 20 random directions, the sham, and five fixed clinical directions: Atelectasis, Pneumothorax, Cardiomegaly, Mass, and Nodule. It produces 28 aggregate configurations and 11,200 per-image outcomes.

For direction $d$, let $\Delta_d$ be its paired mean change from baseline. The familywise statistic is
\begin{equation}
M=\DeltaEff-\max_{d\in\mathcal D_{\mathrm{clinical}}}\Delta_d.
\label{eq:margin}
\end{equation}
The maximum is recomputed inside each of 5,000 patient-bootstrap draws. Direction specificity requires a positive Effusion effect above random p95 and absolute sham, together with a one-sided 95\% lower bound for $M$ above zero.

\paragraph{Evidence and review.}
Every run uses an immutable commit snapshot and records commands, environment, timings, outputs, and terminal checksums. Internal validation replays decision summaries from accepted artifacts. A pinned, read-only cross-family reviewer checks each hard gate after internal verification; review can reject a gate but cannot create experimental evidence.
